\documentclass[11pt]{article}
\usepackage[utf8]{inputenc}
\usepackage[T1]{fontenc}
\usepackage[spanish]{babel}
\usepackage[margin=2.3cm]{geometry}
\usepackage{amsmath,amssymb}
\usepackage{enumitem}
\setlength{\parindent}{0pt}
\setlength{\parskip}{6pt}

\begin{document}

\begin{center}
{\Large\bfseries Control --- Macroeconomía II}\\[4pt]
{\large Mini control: Solow con progreso tecnológico}\\[4pt]
40 puntos \quad$\cdot$\quad 30 minutos \quad$\cdot$\quad sin apuntes \quad$\cdot$\quad con calculadora
\end{center}

\hrule
\medskip

\textit{Sea claro con los gráficos. Nombre los ejes, los puntos de equilibrio y todas las curvas. Diferencie movimientos y desplazamientos de las curvas.}

Considere el modelo de Solow visto en clase con progreso tecnológico exógeno:
\[
Y = K^{\alpha}(AL)^{1-\alpha}, \qquad \frac{\dot A}{A} = g, \qquad \frac{\dot L}{L} = n, \qquad \dot K = sY - \delta K, \qquad 0<\alpha<1,
\]
con $\tilde k \equiv \dfrac{K}{AL}$, $\tilde y \equiv \dfrac{Y}{AL}$ (por trabajador efectivo) y $k \equiv \dfrac{K}{L}$, $y \equiv \dfrac{Y}{L}$ (por trabajador).

\begin{enumerate}[leftmargin=*, itemsep=12pt]

\item \textbf{(15 puntos)}
\begin{enumerate}[label=\alph*), leftmargin=*, itemsep=4pt]
\item \textbf{(3 pts)} Escriba la función de producción en unidades de trabajo efectivo, $\tilde y = f(\tilde k)$.
\item \textbf{(5 pts)} Muestre paso a paso que
\(
\dot{\tilde k} = s\,\tilde k^{\alpha} - (n + g + \delta)\,\tilde k
\).
Explique qué representa cada término de $(n+g+\delta)$.
\item \textbf{(3 pts)} Encuentre $\tilde k^{*}$ e $\tilde y^{*}$ en función de $s$, $n$, $g$, $\delta$ y $\alpha$.
\item \textbf{(4 pts)} Con $\alpha = \tfrac{1}{3}$, $s = 0{,}20$, $n = 0{,}01$, $g = 0{,}02$ y $\delta = 0{,}05$, calcule $\tilde k^{*}$, $\tilde y^{*}$ y $\tilde c^{*}$. Verifique $\tilde c^{*}$ por un segundo camino.
\end{enumerate}

\item \textbf{(10 puntos)} La economía se encuentra en su sendero de crecimiento balanceado cuando, en $t_0$, un terremoto destruye la mitad del stock de capital $K$.
\begin{enumerate}[label=\alph*), leftmargin=*, itemsep=4pt]
\item \textbf{(5 pts)} ¿Qué ocurre en $t_0$ con el nivel de $y$ y con su tasa de crecimiento? ¿Es mayor o menor que $g$? Justifique con $\dot{\tilde k}/\tilde k$.
\item \textbf{(5 pts)} ¿La economía vuelve al mismo sendero de $\ln y$ que antes del terremoto o a uno paralelo? Grafique el diagrama de Solow en $\tilde k$ y $\ln y$ contra el tiempo.
\end{enumerate}

\item \textbf{(5 puntos)} Verdadero o Falso. Según el modelo de Solow visto en clase, en el sendero de crecimiento balanceado la razón capital-producto, $K/Y$, crece a tasa $g$.

\item \textbf{(10 puntos)} Verdadero o Falso. Según el modelo de Solow visto en clase, con mercados competitivos, en el sendero de crecimiento balanceado el salario real crece a tasa $g$, mientras que la tasa de interés real (el producto marginal del capital) permanece constante.

\end{enumerate}

\medskip
\hrule
\smallskip
{\small\textit{Tiempo:} P1 11 min, P2 8 min, P3 4 min, P4 7 min. \textit{Chequeos:} que no se caiga el $\ln$; $\dot{\tilde k}$ escrito completo y su signo a cada lado de $\tilde k^{*}$; «por trabajador efectivo» o «por trabajador», nunca solo «capital»; en V/F, el veredicto es la primera palabra.}

\end{document}
